\documentclass[11pt]{article}

\usepackage[a4paper,left=2.2cm,right=2.2cm,top=1.8cm,bottom=1.8cm]{geometry}
\usepackage[T1]{fontenc}
\usepackage[utf8]{inputenc}
\usepackage{lmodern}
\usepackage{xcolor}
\usepackage{enumitem}
\usepackage[hidelinks]{hyperref}

\setlength{\parindent}{0pt}
\setlength{\parskip}{0.45em}
\setlist[itemize]{leftmargin=*,itemsep=0.1em,topsep=0.2em}

\definecolor{responseblue}{RGB}{0,72,125}

\begin{document}

\begin{center}
{\Large\bfseries Response to the Editor}\\[0.5em]
{\itshape Manuscript: ``An attractor account of numerosity perception''}\\[0.25em]
{\itshape Manuscript ID: IPE-26-0055}
\end{center}

Dear Editor,

Thank you for the opportunity to revise this manuscript. The two reports were careful and technically detailed, and they led us to rebuild the principal analysis and narrow the theoretical claims. We have provided separate point-by-point responses to Reviewer~1 and Reviewer~2.

The original submission was prepared as a Short Report. The requested trial-level analyses and the methodological detail needed to report them transparently take the revision beyond that format. We have therefore prepared the revised manuscript as a Standard Article.

The title has been changed from ``A Hopfield attractor account of numerosity perception'' to ``An attractor account of numerosity perception''. This reflects the revised paper's more precise scope: it tests associative clean-up on a one-dimensional numerical landscape, without invoking the high-dimensional storage-capacity claims usually associated with Hopfield networks.

The principal revisions are as follows:

\begin{itemize}
\item We replaced the summary-curve objective with a categorical trial-level likelihood over all 52,995 integer reports. Response time is now modelled as a separate trial-level log-normal channel.

\item We replaced the original Poisson-well parameterisation with integer-centred Gaussian memory components whose widths increase with magnitude according to a fitted power law.

\item The primary model now has one representational landscape shared across all five exposure durations. Only the sensory channel and lapse vary by duration.

\item We added matched static and constant-width controls, together with a purely driven accumulator, to isolate the contributions of magnitude-dependent broadening, attractor clean-up, and online accumulation. The revised trial-level comparison strongly favours the power-law attractor.

\item Response time is now derived from bounded attractor-convergence time. Convergence alone and a convergence-plus-entropy link fit almost identically, whereas entropy alone performs substantially worse; the simpler convergence-only account is therefore retained.

\item We removed the proposed density regime and its comparison with density-threshold studies. The central claim is now restricted to a shared numerical readout spanning more categorical and more approximate responses.

\item We expanded the treatment of visuospatial-capacity, recurrent visual, image-computable, and other dynamical accounts. The revised manuscript states explicitly that the model begins after scalar numerical evidence has been extracted and does not explain continuous image statistics, segmentation, spatial configuration, or attentional load.

\item We added an all-duration hierarchical individual-differences analysis using disjoint numerosity ranges and a participant lapse effect. Greater fitted landscape resolution is associated with lower Weber noise conditional on lapse; we present this as a correlational, cross-participant result.

\item We removed the developmental-dyscalculia extension.
\end{itemize}

Together, these changes make the response likelihood appropriate to the data, test the dynamics against matched controls, tie the response-time prediction to the dynamics themselves, and align the claims with the explanatory level of the model. We thank you and the reviewers for prompting a substantially stronger and clearer manuscript.

Sincerely,\\
Eudald Correig-Fraga

\end{document}
